\chapter{Conclusions}

This thesis started from a practical gap in 3D Gaussian Splatting: although the method produces photorealistic renderings, the underlying set of Gaussians is optimized for image fidelity rather than for explicit, metrically usable geometry (Sections~\ref{subsec:baseline_limitations}, \ref{sec:results_baseline_compare}). Two complementary strategies were investigated to narrow this gap: depth-supervised optimization using monocular priors from the Depth Anything family, and a structural reparameterization that couples Gaussians to explicit B\'ezier surfaces with object-centric, segmentation-driven initialization. The experiments on two object-centric captures (pizza and parking meter) show that the structural route, when paired with disciplined seed preprocessing and mask supervision, yields meshes that are cleaner and easier to export than the segmented 3DGS baseline.

\section{Main Contributions}

The work makes the following contributions.

\begin{itemize}
    \item \textbf{Depth-supervised 3DGS.} A depth term built on Depth Anything priors was integrated into the training objective (Sections~\ref{subsec:depth_loss}--\ref{subsec:depth_joint_opt}). Qualitatively, depth supervision changes the internal Gaussian layout and depth maps even when final RGB is similar, reducing the depth-stratification artifacts that arise from a purely photometric loss (Section~\ref{subsec:depth_rgb_tradeoff}).
    \item \textbf{B\'ezier structural model coupled to splatting.} Gaussians were reparameterized as samples of explicit cubic B\'ezier surfaces, with means and scales recomputed each forward pass so that gradients flow back to the control points (Sections~\ref{subsec:bezier_representations}--\ref{subsec:bezier_gradients}). This turns an unstructured point cloud into a controllable, exportable surface representation.
    \item \textbf{Object-centric initialization and seed hygiene.} A pipeline combining Grounding DINO and SAM~2 segmentation, COLMAP seed labeling, $k$-NN outlier rejection, and two initialization modes (open PCA patches and a cube AABB lattice) was developed to place control points sensibly before optimization (Sections~\ref{subsec:semantic_isolation}--\ref{subsec:knn_outlier}, \ref{subsec:bezier_init}).
    \item \textbf{Systematic ablation.} A full ablation over initialization geometry, loss formulation ($L_1$/$L_2$, full-frame vs.\ decomposed FgBg mask loss), and seed hygiene was reported with Chamfer distance and control-point gradient norms (Sections~\ref{sec:results_ablation}--\ref{sec:results_comparative}). The study identifies \emph{outlier management} as the single strongest factor: on the pizza scene it lowers unidirectional Chamfer from roughly $21$--$24$ to below $1$ in cube mode, and it visibly stabilizes control-point gradients. Mask-based FgBg supervision is the second key ingredient once seeds are clean, and the \textbf{Cube L1 FgBg Out} configuration gave the best overall trade-off between mesh fidelity, training stability, and visual cleanliness on both scenes.
\end{itemize}

\section{Limitations}

Several limitations should be stated plainly.

\begin{itemize}
    \item \textbf{Fixed topology from initialization.} Both open and cube modes fix the patch layout and, in practice, the surface topology at initialization; training deforms the lattice but cannot change its connectivity. Objects whose shape departs strongly from the initial scaffold, or scenes with multiple disconnected parts, are therefore poorly served by a single AABB-derived lattice (Section~\ref{subsec:bezier_representations}).
    \item \textbf{Dependence on segmentation and seed quality.} The whole pipeline inherits the failure modes of its front end: missed or leaking SAM~2 masks and spurious COLMAP seeds propagate into distorted control nets. Outlier filtering mitigates but does not eliminate this dependence (Section~\ref{subsec:seed_preprocessing}).
    \item \textbf{Limited evaluation scope.} Results cover two object-centric scenes with reference geometry from COLMAP seeds rather than ground-truth scans, and emphasize Chamfer distance and qualitative inspection. Larger, scanned, or full-scene benchmarks were out of scope and would be needed to claim generality.
    \item \textbf{Depth-prior consistency.} Monocular depth priors are only multi-view consistent up to scale and can themselves be wrong on reflective or textureless regions, bounding how much the depth term can help.
\end{itemize}

\section{Future Research Directions}

The most direct extension concerns the \textbf{initialization geometry}. The cube mode in this work anchors control points to the axis-aligned bounding box (AABB) of the segmented seeds (Section~\ref{subsec:bezier_cube}); this is convenient but discards object orientation and forces a six-face layout regardless of shape. A promising direction is to replace the AABB with a learned \textbf{3D box (or instance) segmentation} stage and to fit the B\'ezier scaffold directly to that output. Concretely, an oriented 3D bounding box, or a full 3D instance mask, could be predicted from the point cloud or from multi-view masks using models in the spirit of SAM3D~\cite{zhang2023sam3d} (2D SAM masks lifted to 3D boxes) or native point-cloud segmenters such as Point-SAM~\cite{zhou2025pointsam}. An oriented box immediately yields a better-aligned, shape-aware scaffold than an AABB, while a 3D instance mask would allow the lattice to follow the actual object hull and to support multiple disconnected parts, removing the single-box assumption that currently limits topology. Going one step further, the 3D segmentation could drive the reconstruction end to end: instance masks would both seed the patches and supply per-part mask supervision, tightening the link between the segmentation front end and the structural model.

Several further directions follow naturally from the limitations above.

\begin{itemize}
    \item \textbf{More robust outlier management.} The ablation study showed that $k$-NN rejection on COLMAP seeds is the single strongest lever for both Chamfer distance and control-point gradient stability (Section~\ref{sec:results_comparative}), yet the current filter is a one-shot, distance-based heuristic applied before patch fitting. Future work could treat outlier handling as a first-class, iterative component: multi-stage filtering that combines geometric consistency, reprojection error, and mask agreement across views; confidence-weighted seeding that down-weights rather than discards uncertain points; or learned outlier detectors trained on synthetic mis-segmentation. Tighter coupling with the segmentation front end---for example, rejecting seeds whose 3D neighborhood disagrees with the SAM~2 mask in most views---would reduce the sensitivity to spurious object labels and reflective artifacts that still distort control nets after a single $k$-NN pass.
    \item \textbf{Multi-object scene segmentation and reconstruction.} The reported pipeline is object-centric: Grounding DINO and SAM~2 isolate one target per run, and both test scenes contain a single foreground object. Extending the approach to \emph{all} objects in a scene would require instance segmentation on every view (or a native 3D instance partition of the COLMAP cloud), followed by per-instance seed extraction, outlier filtering, and B\'ezier scaffold fitting. Each object could then carry its own patch lattice and mask-supervision terms while sharing the same camera calibration and rasterizer. Promptable 3D segmenters such as Point-SAM~\cite{zhou2025pointsam}, or multi-view lifting of SAM~2 masks to consistent 3D instances, are natural building blocks for such an extension and would move the method from curated object captures toward cluttered indoor or outdoor scenes relevant to robotics and digital-twin workflows.
    \item \textbf{Adaptive and learnable topology.} Beyond a fixed lattice, patch count and connectivity could change during training (splitting a patch where detail is missing, merging patches that overlap). At patch boundaries, continuity must be enforced: $C^0$ means adjacent patches meet without gaps (as in cube mode, where shared control points on cube edges already guarantee this); $C^1$ additionally aligns tangent directions so the surface does not form visible creases. Future work could learn how to place and connect patches automatically instead of fixing the layout at initialization.
    \item \textbf{Comparison and hybridization with surface-aligned splatting.} Methods such as SuGaR~\cite{guedon2024sugar} extract editable meshes from 3DGS via surface-alignment regularization and Poisson reconstruction. Benchmarking the B\'ezier model against, or hybridizing it with, such surface-aligned approaches would clarify when explicit parametric patches are preferable to post-hoc meshing.
    \item \textbf{Stronger and self-consistent depth supervision.} Multi-view-consistent or metric depth, together with normal priors, could be combined with the structural model to further constrain geometry in weakly observed regions.
    \item \textbf{Broader evaluation.} Validation on scanned ground truth and on full, multi-object scenes, with standardized geometric metrics, is needed to establish how far the approach generalizes beyond the two test objects.
\end{itemize}

Taken together, these directions point toward a reconstruction pipeline in which \emph{semantic 3D segmentation}---extended to every object in the scene, paired with stronger outlier control---drives explicit and editable surface models, while differentiable splatting continues to supply the photometric signal. Such a pipeline would be directly useful in the application settings that motivated this thesis, including robotics, augmented reality, digital twins, and simulation, where compact, metrically meaningful, and editable geometry matters as much as photorealism.
